\documentclass[12pt]{article}
\usepackage[utf8]{inputenc}
\usepackage[T1]{fontenc}
\usepackage[margin=1in]{geometry}
\usepackage{setspace}
\usepackage{xcolor}
\usepackage[colorlinks=true, urlcolor=blue!60!black, linkcolor=blue!60!black]{hyperref}
\setlength{\parindent}{0em}
\setlength{\parskip}{0.6em}
\renewcommand{\baselinestretch}{1.0}
\begin{document}
\begin{center}
{\LARGE \textit{Counterman v. Colorado}, 600 U.S. 66 (2023)\par}
\vspace{0.5em}
{\large Case Brief --- Two-Page Edition\par}
\end{center}

\section{Memory Jogger}
A true-threats prosecution requires proof that the speaker was at least reckless about whether his words would be seen as threatening violence, not merely that a reasonable person would see a threat (\href{https://supreme.justia.com/cases/federal/us/600/22-138/}{600 U.S. 66}).

\section{Detailed Case Facts}
Billy Counterman sent hundreds of unsolicited Facebook messages to singer Coles Whalen over two years, opening new accounts when blocked; some said ``Was that you in the white Jeep?'' and ``Staying in cyber life is going to kill you.'' Terrified, Whalen stopped walking alone and canceled performances.

\section{Procedural History}
Colorado convicted Counterman of stalking under an objective true-threats test; the Colorado Court of Appeals affirmed, and the Supreme Court granted certiorari.

\section{Judicial Votes}
Vacated and remanded, 7--2. Kagan wrote for Roberts, Alito, Kavanaugh, and Jackson; Sotomayor, joined in part by Gorsuch, concurred in part and in the judgment; Thomas and Barrett dissented.

\section{Holding}
The First Amendment requires the State to prove at least recklessness: that the defendant ``consciously disregarded a substantial risk that his communications would be viewed as threatening violence.''

\section{Analysis of Opinions}
Kagan held that although threats are judged objectively, the First Amendment gives ``strategic protection'' to avoid chilling protected speech, as with actual malice in defamation, intent in incitement, and scienter in obscenity. Recklessness balances victims' interests against chill, matches \textit{Sullivan}, and stops short of the intent required for incitement, which sits ``a hair's-breadth away'' from political advocacy. Sotomayor agreed for stalking but would require purpose or knowledge for pure threats. Thomas attacked reliance on \textit{Sullivan}, and Barrett defended an objective standard, calling recklessness a ``Goldilocks judgment'' that hampers protective orders and discipline.

\section{Examples: Future Applications}
\textbf{Same-Side:} a conviction for a heated livestream comment under a ``reasonable viewer'' instruction is invalid. \textbf{Opposite-Side:} 150 ominous messages sent after the recipient wrote ``Your messages terrify me'' support recklessness. \textbf{Fence-Sitter:} a civil protective order issued under an objective standard.

\section{Critique}
Barrett's delusional-stalker paradox: the most dangerous stalkers may lack awareness that their words frighten. The majority's open interest-balancing drew fire from both originalists and speech-protective Justices.

\section{Key Quotations}
\begin{quote}
``We hold that it does, but that a mental state of recklessness is sufficient. The State must show that the defendant consciously disregarded a substantial risk that his communications would be viewed as threatening violence. The State need not prove any more demanding form of subjective intent to threaten another.''
\end{quote}
\noindent --- Justice Kagan, Opinion of the Court, 600 U.S. at 69.

\end{document}
